% Engine setup, palette and link handling.
%
% altacv.cls is vendored verbatim from upstream v1.1.5, so every deviation
% from stock AltaCV lives in this file rather than in the class.

\usepackage[english]{babel}

% AltaCV puts the sidebar in the right margin, so the text block is narrow and
% the margin is wide. marginparwidth + marginparsep must stay inside `right`
% or the sidebar runs off the page edge.
\geometry{
  left           = 2cm,
  right          = 10cm,
  marginparwidth = 6.8cm,
  marginparsep   = 1.2cm,
  top            = 1.25cm,
  bottom         = 1.25cm,
}

% Tectonic drives XeTeX, and fontspec resolves names against *installed system
% fonts*. Lato is not installed on a CI runner, so it is loaded by file name
% instead: those .ttf files ship inside the TeX Live bundle that tectonic
% downloads, which is what makes a laptop build and a runner build identical.
% Lato is also what upstream's pdfLaTeX branch selects, so the typeface is
% unchanged from the original template.
\ifxetexorluatex
  \setmainfont{Lato}[
    Extension      = .ttf,
    UprightFont    = *-Regular,
    BoldFont       = *-Bold,
    ItalicFont     = *-Italic,
    BoldItalicFont = *-BoldItalic,
  ]
\else
  \usepackage[utf8]{inputenc}
  \usepackage[T1]{fontenc}
  \usepackage[default]{lato}
\fi

% ragged2e caps the right-hand slack at 2em, which is too tight for the dense
% technology lists: two of them overran the text block by 5pt and 10pt rather
% than wrapping one word earlier. More slack removes both overfull boxes.
\setlength{\RaggedRightRightskip}{0pt plus 3em}

% The class sets itemsep=0.25\baselineskip. Across the ~30 bullets in this
% document that padding alone costs most of a page, and the bullets are
% single ideas that do not need the separation to stay legible. Tightening it
% is what keeps the CV to two pages without deleting anything a reader wants.
\setlist{itemsep=0.12\baselineskip}

% Monochrome palette. The class expects these four colour aliases; the names
% below describe what they actually are, unlike the upstream template which
% defined a "VividPurple" that was set to pure black.
\definecolor{Ink}{HTML}{000000}
\definecolor{Graphite}{HTML}{2E2E2E}
\colorlet{heading}{Ink}
\colorlet{accent}{Ink}
\colorlet{emphasis}{Graphite}
\colorlet{body}{Graphite}

\renewcommand{\itemmarker}{{\small\textbullet}}
\renewcommand{\ratingmarker}{\faCircle}

% The email glyph is an envelope rather than the stock \faAt. FontAwesome 4
% names the "at" glyph `_475`, which xdvipdfmx cannot parse into a ToUnicode
% entry, and it was the single source of the build's only warning -- every
% other icon in this CV is clean. An envelope also matches the weight of the
% phone, marker, LinkedIn and GitHub glyphs beside it; \faAt sat among them
% as the odd one out. For a literal @ instead, use {@} here.
\renewcommand{\emailsymbol}{\faEnvelope}

% This PDF is read in a browser tab, so it carries real document metadata:
% the title is what the tab and the download dialog show.
\usepackage[hidelinks]{hyperref}
\hypersetup{
  pdftitle   = {Oleksandr Ponomarov - Curriculum Vitae},
  pdfauthor  = {Oleksandr Ponomarov},
  pdfsubject = {Senior DevOps Engineer},
  pdflang    = {en},
}

% url.sty defaults \url to \ttfamily, and nothing else in this CV is
% monospaced: the one footnote URL was being set in Latin Modern Mono against
% a Lato page. `same` keeps it in the surrounding typeface while url.sty
% still handles the line breaking.
\urlstyle{same}

% Deliberately no page numbers, even at three pages. The class sets
% \pagestyle{empty}, and adding a footer here needs more than a pagestyle:
% geometry's bottom=1.25cm leaves no band for one (the footer annotation
% lands below the MediaBox), and ragged2e's finite \rightskip makes the
% footer box underfull on every shipout. Both are fixable, but only by
% retuning the page geometry this layout depends on.

% Same glyphs and same text as the stock macros, but each contact detail is
% now actionable for a reader who opens the CV on a screen. AltaCV v1.1.5
% predates upstream's `withhyper` option, so the five macros are wrapped here.
\renewcommand{\email}[1]{\printinfo{\emailsymbol}{\href{mailto:#1}{#1}}}
\renewcommand{\phone}[1]{\printinfo{\phonesymbol}{\href{tel:#1}{#1}}}
\renewcommand{\homepage}[1]{\printinfo{\homepagesymbol}{\href{https://#1}{#1}}}
\renewcommand{\linkedin}[1]{%
  \printinfo{\linkedinsymbol}{\href{https://www.linkedin.com/in/#1}{#1}}}
\renewcommand{\github}[1]{%
  \printinfo{\githubsymbol}{\href{https://github.com/#1}{#1}}}
